% ATTEMPT_STATUS: unresolved
% RESEARCH_GENERATED: 2026-10-01; GPT-6 Astra Ultra; individual mathematical attempt
% TOP_PROBLEM: {"id": 103, "title": "4-term APs in Fourier Uniform Sets", "category_id": 2, "category": {"id": 2, "name": "combinatorics", "display_name": "Combinatorics", "description": "Counting problems, graph theory, discrete structures.", "slug": "combinatorics", "order_index": 2, "created_at": "2026-07-31T15:26:25.671Z"}, "status": "open"}
% FIRST_POSTED: 2026-10-01
% Imported from: unsolvedmath_additional_500.tex; UnsolvedMath ID 103
% !TeX program = lualatex
% Compile twice: lualatex 103.tex
% Self-contained: no external bibliography, images, or \input files.
% Numeric section labels and visible numbers are original UnsolvedMath IDs.
\documentclass[11pt,a4paper]{article}
\usepackage[margin=24mm,headheight=14pt]{geometry}
\usepackage{fontspec}
\setmainfont{Latin Modern Roman}
\setsansfont{Latin Modern Sans}
\setmonofont{Latin Modern Mono}
\usepackage{amsmath,amssymb,mathtools}
\usepackage{unicode-math}
\setmathfont{Latin Modern Math}
\usepackage{microtype}
\usepackage{enumitem,xurl,needspace}
\usepackage[dvipsnames]{xcolor}
\usepackage{hyperref}
\hypersetup{unicode=true,colorlinks=true,linkcolor=MidnightBlue,urlcolor=MidnightBlue,
 pdftitle={UnsolvedMath Problem 103},pdfauthor={Source-annotated compilation},bookmarksnumbered=true}
\usepackage{bookmark,tocloft,titlesec,fancyhdr}
\setlength{\cftsecnumwidth}{4.2em}
\cftsetpnumwidth{2.5em}
\cftsetrmarg{4em}
\titleformat{\section}{\Large\bfseries\raggedright}{\thesection}{0.7em}{}
\pagestyle{fancy}\fancyhf{}
\fancyhead[L]{\small\sffamily UnsolvedMath research attempt}
\fancyhead[R]{\small Original catalogue IDs}
\fancyfoot[C]{\thepage}
\setlength{\parindent}{0pt}
\setlength{\parskip}{5pt plus 1pt}
\setlength{\emergencystretch}{3em}
\setcounter{tocdepth}{1}\setcounter{secnumdepth}{1}
\setlist[itemize]{leftmargin=3em,itemsep=4pt,topsep=4pt,parsep=0pt}
\allowdisplaybreaks
\newcommand{\N}{\mathbb{N}}
\newcommand{\Z}{\mathbb{Z}}
\newcommand{\Q}{\mathbb{Q}}
\newcommand{\R}{\mathbb{R}}
\newcommand{\C}{\mathbb{C}}
\newcommand{\F}{\mathbb{F}}
\newcommand{\A}{\mathbb{A}}
\newcommand{\PP}{\mathbb{P}}
\newcommand{\E}{\mathbb{E}}
\newcommand{\eps}{\varepsilon}
\newcommand{\dd}{\,\mathrm d}
\newcommand{\sref}[2]{\hyperlink{source:#1:#2}{[S#2]}}
\newif\ifshowcatalogue
\showcataloguetrue
\newcommand{\cataloguescope}[1]{\ifshowcatalogue
 \paragraph{Statement recorded in the supplied source.}
 #1\fi}
\begin{document}
% UnsolvedMath ID 103; source catalogue code GREEN-013

\setcounter{section}{102}

\section{Four-Term Progressions under Fourier Uniformity}\label{103}

{\small\sffamily UnsolvedMath ID 103 · Combinatorics}

\subsection{Definitions and mathematical statement}

\paragraph{Shared notation.}Unless a section says otherwise, $\N=\{1,2,\ldots\}$,
$\Z$ is the ring of integers, $\Q,\R,\C$ are the rational, real, and complex
fields, $[N]=\{1,\ldots,\lfloor N\rfloor\}$, and $\log$ is the natural
logarithm. For real functions of a parameter tending to infinity,
$f\ll g$ and $f=O(g)$ mean $|f|\leq Cg$ eventually for some fixed $C$;
$f\gg g$ reverses this comparison; $f=o(g)$ means $f/g\to0$;
$f\sim g$ means $f/g\to1$; and $f\asymp g$ means both order bounds.
A subscript on $O$ or $\ll$ specifies allowed constant dependence.
The expression $n^{o(1)}$ in an upper bound means growth smaller than
$n^\epsilon$ for each fixed $\epsilon>0$. Local definitions govern any
exception, finite-field convention, or use of the same symbol for another
object. An asymptotic claim is not automatically an effective algorithm.



For $f:\Z/N\Z\to\C$ use the unnormalised transform $\widehat f(r)=\sum_x f(x)e^{-2\pi i rx/N}$. Fourier uniformity means $\max_r|\widehat{1_A-\alpha}(r)|=o(N)$ along the family under consideration. Progressions are counted by $T_4(A)=\sum_{x,d}\prod_{j=0}^3 1_A(x+jd)$. For fixed positive density, degenerate choices contribute only lower-order terms. The question asks for a universal lower-bound constant at the specified power of $\alpha$. \sref{103}{1} \sref{103}{2}

\paragraph{Statement recorded in the supplied source.}
Suppose that $A \subset \mathbb{Z}/N\mathbb{Z}$ has density $\alpha$ and is Fourier uniform (all Fourier coefficients of $1_A - \alpha$ are $o(N)$). Does $A$ contain at least $\gg \alpha^{100}N^2$ 4-term arithmetic progressions? \sref{103}{1}

\paragraph{Residual task reported by the archive.}

The embedded review (2026-08-17) records: Construct Fourier-uniform counterexamples below alpha\textasciicircum{}100, or prove the bound. \sref{103}{1}

\subsection{Short English statement}

If a set looks uniform to every Fourier frequency, must it still contain a polynomial-in-density proportion of all four-term arithmetic progressions?

\subsection{Research attempt}

\paragraph{Provenance and scope.}
This individual attempt was generated by GPT-6 Astra Ultra on 1 October 2026. The method was to derive explicit partial results and test the first proposed extension adversarially. The original statement and sources below are retained. No claim of mathematical novelty or complete solution is made.

\paragraph{Formal quantifiers and primary-source context.}
Use the normalised Fourier transform $\widetilde f(r)=N^{-1}\sum_x f(x)e^{-2\pi irx/N}$. The uniformity assumption is $\max_r|\widetilde{1_A-\alpha}(r)|=o(1)$ along a family, with the density parameter fixed before the ambient size tends to infinity. Deng, Tidor and Zhao's inspected paper gives conditional constructions through a colouring question; its abstract does not assert an unconditional resolution of this four-term problem. We test a frequently tempting proof route: replacing Fourier uniformity by the assertion that every multilinear four-term correlation behaves randomly.

\paragraph{The exact Fourier obstruction.}
For any real $f$ on $\Z/N\Z$, Fourier expansion and orthogonality give
\[
 \E_{x,d}\prod_{j=0}^3f(x+jd)
 =\sum_{r,s}\widetilde f(r+2s)\widetilde f(-2r-3s)
                 \widetilde f(r)\widetilde f(s).
\]
Indeed, if the four frequencies are $a,b,c,d'$, averaging $x$ imposes $a+b+c+d'=0$, and averaging the common difference imposes $b+2c+3d'=0$. Taking $r=c,s=d'$ gives the formula, with no division and therefore no restriction on $N$. The zero-frequency term is $\alpha^4$. The other terms are not absolute squares, so their signs cannot be ignored. Bounding each nonzero Fourier coefficient individually and then summing all frequency pairs loses powers of $N$; the stated $o(1)$ hypothesis has no prescribed rate with which to absorb that loss.

For comparison, three-term counts have only one free frequency. Over an odd prime modulus, Parseval and Cauchy--Schwarz bound the nonzero contribution by
\[
 \sup_{r\ne0}|\widetilde f(r)|
 \Big(\sum_r|\widetilde f(r)|^2\Big)^{1/2}
 \Big(\sum_r|\widetilde f(-2r)|^2\Big)^{1/2}
 \leq \alpha\sup_{r\ne0}|\widetilde f(r)|
\]
for an indicator of density $\alpha$. The permutation $r\mapsto-2r$ is the reason for the odd-modulus restriction in this comparison. There is no analogous one-dimensional summation in the displayed four-term formula.

\paragraph{An explicit mixed-correlation counterexample to the proposed estimate.}
Let $p>3$ be prime and put $u_j(x)=e_p(c_jx^2)$, where
\[
 (c_0,c_1,c_2,c_3)=(1,-3,3,-1),\qquad e_p(t)=e^{2\pi it/p}.
\]
For each $j$ every normalised Fourier coefficient has modulus $p^{-1/2}$. Here is a direct proof of the needed Gauss-sum fact. For $c\ne0$, write $S=\sum_x e_p(cx^2-rx)$. On putting $x=y+h$ in $|S|^2$, the inner sum over $y$ is $\sum_y e_p(2chy)$ times a unit complex factor. It vanishes unless $h=0$, since $2c$ is invertible. The remaining term is $p$, proving $|S|=\sqrt p$.

Nevertheless the mixed four-term correlation is identically one:
\[
 \E_{x,d}\prod_{j=0}^3u_j(x+jd)=1.
\]
Expanding the quadratic polynomial shows why: the coefficients of $x^2$, $xd$, and $d^2$ are respectively $\sum c_j$, $2\sum jc_j$, and $\sum j^2c_j$, all zero. This explicitly refutes any general claim that small ordinary Fourier coefficients force all four-function progression correlations to be small. Subtracting the means changes the average by $O(p^{-1/2})$, because each mean is of that size and all original functions are bounded by one. Thus centring the functions does not remove the obstruction.

\paragraph{What the counterexample does and does not establish.}
The four functions above are different and complex-valued. They are not four copies of one indicator. Therefore their large correlation does not refute the catalogue's positive lower bound. Passing from phases to a set requires controlling level sets and their joint distribution; taking real parts or thresholding can change both the density and the correlation. This is exactly where an attempted quick disproof would become invalid. The construction instead explains, by an explicit calculation, why linear Fourier uniformity leaves quadratic structure available.

\paragraph{Verification and remaining gap.}
The cancellation identities hold over the integers before reduction modulo $p$, and the Gauss-sum calculation explicitly excludes characteristics two and three. Degenerate progressions have density $1/p$ and cannot account for the limiting mixed correlation. The finite-field example is used only as an obstruction to a proof method, not silently substituted for a cyclic-set solution. The unresolved task is a same-indicator lower bound with exponent one hundred, or a valid Fourier-uniform indicator construction violating it. The exact Fourier formula and the quadratic mixed example are the checked output of this attempt; neither supplies that final step.

\subsection{Sources}

{\small

\begin{itemize}

\item[\hypertarget{source:103:1}{S1}] ULAM.AI, \emph{UnsolvedMath}, supplied \texttt{problems.json}, record 103 (GREEN-013), statement and background. The embedded literature-review date is 2026-08-17; that is the archive’s review date, not a fresh certification. Dataset landing page: \url{https://huggingface.co/datasets/ulamai/UnsolvedMath}.

\item[\hypertarget{source:103:2}{S2}] Ben Green, 100 Open Problems (author’s maintained problem list). \url{https://people.maths.ox.ac.uk/greenbj/papers/open-problems.pdf}. The author-hosted document was inspected on 27 September 2026; the archive’s local code is not a problem-number locator in this paper.

\item[\hypertarget{source:103:3}{S3}] M. Deng, J. Tidor, Y. Zhao, Uniform sets with few progressions via colourings, Math. Proc. Camb. Phil. Soc. 179 (2025), 79--103. \url{https://doi.org/10.1017/S0305004125000106}. Bibliographic information imported from the supplied record.

\item[E1] M. Deng, J. Tidor and Y. Zhao, \emph{Uniform sets with few progressions via colorings}, current manuscript. \url{https://arxiv.org/abs/2307.06914}. Inspected 1 October 2026.

\end{itemize}

}

\label{end:103}
\end{document}
